% answers_short.tex —— 1952 年·全国卷 · 简短答案速查
% 由 tools/make_short_answers.py 自 answers.tex 的 \gksolhead/\gkanswer 提取，
% 并与 1952/1952_数学真题_答案要点.md 逐题比对后生成，勿手改（改源 MD 后重跑）。

\begin{gknotice}
\textbf{使用说明}：本节只给最终答案便于快速核对，完整推导见解析版。答案为编写者独立推导，\textbf{非官方标准答案}。
\end{gknotice}

\gkpart{填空题}
\begin{itemize}[leftmargin=2.6em]
  \item[\textbf{1.}] $(x-y)(x+y)(x^2+y^2)$
  \item[\textbf{2.}] $2$
  \item[\textbf{3.}] $-1$
  \item[\textbf{4.}] $\pm3$
  \item[\textbf{5.}] $-24$
  \item[\textbf{6.}] $4\sqrt3$ 寸
  \item[\textbf{7.}] $15$ \mbox{平方寸}
  \item[\textbf{8.}] $144^\circ$
  \item[\textbf{9.}] $\dfrac{22}{7}$，$\dfrac{355}{113}$，$3.14159265$
  \item[\textbf{10.}] $4$ 倍
  \item[\textbf{11.}] $12\pi$ 立方尺
  \item[\textbf{12.}] $5$ 种
  \item[\textbf{13.}] $\dfrac79$
  \item[\textbf{14.}] $\dfrac12\left(n\pi+\dfrac\pi4\right)$（$n$ 为整数）
  \item[\textbf{15.}] $\dfrac{5\sqrt3}{3}$ 丈
  \item[\textbf{16.}] $3$ \mbox{平方寸}
  \item[\textbf{17.}] $x+y+1=0$
  \item[\textbf{18.}] $x^2+y^2-6x-8y=0$
  \item[\textbf{19.}] $\dfrac15$
  \item[\textbf{20.}] $(1,-3)$
\end{itemize}

\gkpart{解答题}
\begin{itemize}[leftmargin=2.6em]
  \item[\textbf{21.}] $2$，$-6$，$\omega$，$\omega^2$（其中 $\omega=\dfrac{-1+\mathrm i\sqrt3}{2}$）
  \item[\textbf{22.}] 见解析版
  \item[\textbf{23.}] (1) $\cos C=\dfrac18$；(2) $\sin C=\dfrac{3\sqrt7}{8}$，$\sin B=\dfrac{5\sqrt7}{16}$，$\sin A=\dfrac{\sqrt7}{4}$；(3) $A$，$B$，$C$ 皆为锐角
  \item[\textbf{24.}] 长轴 $\dfrac{2\sqrt{165}}{3}$，短轴 $\dfrac{2\sqrt{385}}{7}$，焦点 $\left(0,\pm\dfrac{2\sqrt{1155}}{21}\right)$
\end{itemize}
